% Part: lambda-calculus
% Chapter: introduction
% Section: fixed-point-combinator

\documentclass[../../../include/open-logic-section]{subfiles}

\begin{document}

\olfileid{lam}{int}{fix}
\olsection{નિશ્ચિત-બિંદુ સંયોજકો}

ધારો કે આપણી પાસે લૅમ્બડા પદ~$g$ છે અને એવું બીજું પદ~$k$ જોઈએ છે
કે~$k$, $gk$ સાથે $\beta$-સમકક્ષ હોય. આપણા સંકેતન પ્રમાણે પદો
વ્યાખ્યાયિત કરો:
\[
\fn{diag}(x) = xx
\]
અને
\[
l(x) = g(\fn{diag}(x))
\]
બીજા શબ્દોમાં, $l$ એ પદ~$\lambd[x][g(xx)]$ છે. $k$ એ પદ~$ll$
રાખો. ત્યારે મળે છે:
\begin{align*}
k & = (\lambd[x][g(xx)])(\lambd[x][g(xx)]) \\
& \red  g((\lambd[x][g(xx)])(\lambd[x][g(xx)])) \\
& = gk.
\end{align*}
જો
\[
Y = \lambd[g][((\lambd[x][g(xx)])(\lambd[x][g(xx)]))]
\]
રાખીએ, તો $Yg$ અને $g(Yg)$ બંને એક જ પદ સુધી લઘુકરણ પામે છે;
આથી $Yg \equiv_\beta g(Yg)$. તેને ``કરીનો સંયોજક'' કહે છે. તેના
બદલે જો
\[
Y = (\lambd[xg][g(xxg)])(\lambd[xg][g(xxg)])
\]
રાખીએ, તો ખરેખર $Yg$નું $g(Yg)$ સુધી લઘુકરણ થાય છે; આ વધુ
મજબૂત વિધાન છે. $Y$ના આ બીજા રૂપને ``ટ્યુરિંગનો સંયોજક'' કહે છે.

\end{document}
